% !TEX program = lualatex
%=====================================================================
%  Criptografía y Seguridad - ITBA
%  Clase 3 — MACs y Cifrado Autenticado
%  Apunte integral: teoría + práctica
%=====================================================================
\documentclass[11pt,a4paper]{article}

%---------------------------------------------------------------------
%  Paquetes
%---------------------------------------------------------------------
\usepackage{iftex}
\ifPDFTeX
  \usepackage[utf8]{inputenc}
  \usepackage[T1]{fontenc}
  \usepackage{lmodern}
\else
  \usepackage{fontspec}
\fi
\usepackage[spanish,es-tabla,es-noshorthands]{babel}
\usepackage{microtype}

\usepackage{amsmath,amssymb,amsthm}
\usepackage{mathtools}
\usepackage{geometry}
\geometry{a4paper,margin=2.4cm,headheight=22pt}

\usepackage{xcolor}
\usepackage{graphicx}
\usepackage{enumitem}
\usepackage{booktabs}
\usepackage{array}
\usepackage{tcolorbox}
\tcbuselibrary{skins,breakable}
\usepackage{fancyhdr}
\usepackage{titlesec}
\usepackage{hyperref}
\usepackage{listings}
\usepackage{caption}

\usepackage{tikz}
\usetikzlibrary{shapes.geometric,shapes.misc,shapes.symbols,arrows.meta,positioning,calc,
                fit,backgrounds,decorations.pathmorphing,shadows.blur,chains,matrix}

%---------------------------------------------------------------------
%  Paleta de color (inspirada en las diapositivas de la cátedra)
%---------------------------------------------------------------------
\definecolor{itbaCyan}{HTML}{6EC6D9}
\definecolor{itbaCyanDark}{HTML}{2E8B9E}
\definecolor{itbaBlue}{HTML}{AEDCEA}
\definecolor{boxBlue}{HTML}{BBD4F0}
\definecolor{slideGray}{HTML}{ECECEC}
\definecolor{practGreen}{HTML}{4F8A3D}
\definecolor{practGreenL}{HTML}{E2EFD9}
\definecolor{attackRed}{HTML}{C0392B}
\definecolor{codeBg}{HTML}{F5F5F5}

%---------------------------------------------------------------------
%  Hipervínculos
%---------------------------------------------------------------------
\hypersetup{
  colorlinks=true,
  linkcolor=itbaCyanDark,
  urlcolor=itbaCyanDark,
  citecolor=itbaCyanDark,
  pdftitle={Clase 3 - MACs y Cifrado Autenticado},
  pdfauthor={Criptografia y Seguridad - ITBA}
}

%---------------------------------------------------------------------
%  Encabezados y pies
%---------------------------------------------------------------------
\pagestyle{fancy}
\fancyhf{}
\lhead{\small\textsc{Criptografía y Seguridad — ITBA}}
\rhead{\small Clase 3 · MACs y Cifrado Autenticado}
\cfoot{\small\thepage}
\renewcommand{\headrulewidth}{0.4pt}
\renewcommand{\headrule}{\hbox to\headwidth{\color{itbaCyanDark}\leaders\hrule height \headrulewidth\hfill}}

%---------------------------------------------------------------------
%  Títulos de sección con barra de color (estilo diapositiva)
%---------------------------------------------------------------------
\titleformat{\section}
  {\normalfont\Large\bfseries\color{itbaCyanDark}}
  {\thesection}{0.6em}{}
  [\vspace{-0.4em}{\color{itbaCyan}\titlerule[2pt]}]
\titleformat{\subsection}
  {\normalfont\large\bfseries\color{itbaCyanDark}}
  {\thesubsection}{0.6em}{}
\titleformat{\subsubsection}
  {\normalfont\normalsize\bfseries\color{black}}
  {\thesubsubsection}{0.6em}{}

%---------------------------------------------------------------------
%  Cajas reutilizables
%---------------------------------------------------------------------
% Caja "diapositiva" (recuadro gris de las slides)
\newtcolorbox{slidebox}[1][]{
  enhanced, colback=slideGray, colframe=black!55, boxrule=0.5pt,
  arc=1pt, left=6pt,right=6pt,top=4pt,bottom=4pt, #1}

% Caja de definición
\newtcolorbox{defbox}[1]{
  enhanced, breakable, colback=itbaBlue!25, colframe=itbaCyanDark,
  boxrule=1pt, arc=2pt, left=8pt,right=8pt,top=6pt,bottom=6pt,
  fonttitle=\bfseries, coltitle=white, title={#1}}

% Caja de nota práctica
\newtcolorbox{practbox}[1]{
  enhanced, breakable, colback=practGreenL, colframe=practGreen,
  boxrule=1pt, arc=2pt, left=8pt,right=8pt,top=6pt,bottom=6pt,
  fonttitle=\bfseries, coltitle=white, title={#1}}

% Caja de atención / ataque
\newtcolorbox{warnbox}[1]{
  enhanced, breakable, colback=attackRed!8, colframe=attackRed,
  boxrule=1pt, arc=2pt, left=8pt,right=8pt,top=6pt,bottom=6pt,
  fonttitle=\bfseries, coltitle=white, title={#1}}

%---------------------------------------------------------------------
%  Estilo de código
%---------------------------------------------------------------------
\lstdefinestyle{plain}{
  basicstyle=\ttfamily\footnotesize,
  backgroundcolor=\color{codeBg},
  keywordstyle=\color{itbaCyanDark}\bfseries,
  commentstyle=\color{practGreen}\itshape,
  stringstyle=\color{attackRed},
  breaklines=true, frame=single, rulecolor=\color{black!30},
  showstringspaces=false, columns=fullflexible, tabsize=2,
  xleftmargin=10pt, framexleftmargin=8pt}

%---------------------------------------------------------------------
%  Macros matemáticos
%---------------------------------------------------------------------
\newcommand{\Kset}{\mathcal{K}}
\newcommand{\Pset}{\mathcal{P}}
\newcommand{\Cset}{\mathcal{C}}
\newcommand{\Tset}{\mathcal{T}}
\newcommand{\Sset}{\mathcal{S}}
\newcommand{\xor}{\oplus}
\newcommand{\concat}{\,\|\,}
\newcommand{\Mac}{\mathsf{Mac}}
\newcommand{\Vrfy}{\mathsf{Vrfy}}
\newcommand{\Gen}{\mathsf{Gen}}
\newcommand{\Enc}{\mathsf{Enc}}
\newcommand{\Dec}{\mathsf{Dec}}
\DeclareMathOperator{\negl}{negl}
\DeclareMathOperator{\Trunc}{Truncate}

%=====================================================================
\begin{document}
%=====================================================================

%---------------------------------------------------------------------
%  Portada
%---------------------------------------------------------------------
\begin{titlepage}
\centering
\vspace*{1.2cm}

\begin{tikzpicture}
  \node[rounded corners=2pt, fill=itbaCyan, text=black, inner sep=14pt,
        minimum width=0.8\textwidth, align=center] (t)
        {\Huge\bfseries Criptografía y Seguridad};
\end{tikzpicture}

\vspace{0.6cm}
{\LARGE Integridad y autenticación de mensajes}\\[0.25cm]
{\Huge\bfseries\color{itbaCyanDark} MACs y Cifrado Autenticado}

\vspace{0.8cm}

% Ilustración: bóveda / caja fuerte (recreación del icono de portada)
\begin{tikzpicture}[scale=1.0]
  \fill[black!12] (-3.2,-2.4) rectangle (3.2,2.8);
  \draw[line width=2pt, black!55] (-3.2,-2.4) rectangle (3.2,2.8);
  \fill[itbaCyanDark!75] (-2.4,-2.0) rectangle (2.4,2.4);
  \draw[line width=3pt, black!70] (-2.4,-2.0) rectangle (2.4,2.4);
  \draw[line width=1.2pt, itbaCyan] (-2.0,-1.6) rectangle (2.0,2.0);
  \draw[line width=2pt, black!70] (0,-2.0) -- (0,2.4);
  \foreach \a in {0,45,...,315}{
    \draw[line width=2pt, black!70] (-1.1,0.2) -- ++(\a:0.6);
  }
  \draw[line width=2.5pt, black!70] (-1.1,0.2) circle (0.42);
  \fill[itbaCyan] (-1.1,0.2) circle (0.16);
  \draw[line width=2.5pt, black!70] (1.1,0.2) circle (0.42);
  \fill[itbaCyan] (1.1,0.2) circle (0.16);
  \foreach \y in {-1.2,-0.4,0.8,1.6}{
     \fill[black!70] (-2.2,\y) circle (0.07);
     \fill[black!70] (2.2,\y) circle (0.07);
  }
\end{tikzpicture}

\vfill
{\large Apunte integral — Teoría + Práctica}\\[0.2cm]
{\normalsize Clase 3 \quad·\quad Capítulo 4, \emph{Introduction to Modern Cryptography} (Katz \& Lindell)}\\[0.4cm]
{\small Instituto Tecnológico de Buenos Aires (ITBA)}

\vspace{0.6cm}
{\footnotesize\itshape Documento elaborado a partir de las transcripciones de la clase
teórica (Prof.\ Pablo Abad) y de la clase práctica, junto con las diapositivas oficiales de ambas.}
\end{titlepage}

\tableofcontents
\newpage

%=====================================================================
\section{Motivación: del cifrado a la integridad}
%=====================================================================

Hasta esta clase la criptografía vista resolvía un único gran problema:
\textbf{la confidencialidad} (cómo controlar la diseminación de la información). Esta clase
abre un \textbf{nuevo campo}: el de la \textbf{integridad}, es decir, \emph{cómo confiar en
que la información no fue modificada}, aun cuando provenga de una fuente confiable.

\begin{slidebox}
\textbf{Repaso — Criptosistema.} Es una terna de algoritmos
$\langle \Gen, \Enc, \Dec\rangle$:
\[
  \Gen:()\to\Kset,\qquad \Enc:\Kset\times\Pset\to\Cset,\qquad \Dec:\Kset\times\Cset\to\Pset
\]
con la \textbf{propiedad básica} $\;d_k(e_k(m))=m\;$ para todo $m$ y $k$ válidos.
Conjuntos: $\Kset$ (claves), $\Pset$ (planos), $\Cset$ (cifrados).
\end{slidebox}

\subsection{Recordatorio: el escenario CPA}
El mejor escenario de seguridad visto hasta ahora era el de \textbf{ataque de texto plano
escogido} (\emph{Chosen Plaintext Attack}). En su experimento de indistinguibilidad
$\mathrm{CPA}_{A,\Pi}$, el adversario $A$ tiene acceso a un \textbf{oráculo de cifrado}
$f(x)=e_k(x)$ (nunca conoce la clave):

\begin{defbox}{Experimento $\mathrm{CPA}_{A,\Pi}$}
\begin{enumerate}[leftmargin=1.6em,itemsep=1pt]
  \item Se genera una clave $k\leftarrow \Kset$.
  \item $A$ obtiene $f(x)=e_k(x)$ y genera $(m_0,m_1)$.
  \item Se genera $b\leftarrow\{0,1\}$.
  \item $A$ recibe $c=e_k(m_b)$.
  \item $A$ emite $b'\in\{0,1\}$.
\end{enumerate}
$\mathrm{CPA}_{A,\Pi}=1$ si $b=b'$ (A gana). Si $\Pr[\mathrm{CPA}_{A,\Pi}=1]=\tfrac12+\varepsilon$
con $\varepsilon$ despreciable, $\Pi$ es \textbf{CPA-secure} (indistinguible).
\end{defbox}

Los criptosistemas \textbf{CPA-secure} son no deterministas: incorporan un parámetro extra
público y aleatorio, el \textbf{IV} (\emph{tweak}), que evita reutilizar la misma semilla del
generador. Son CPA-secure los cifrados de flujo con IV y los de bloque con encadenamiento
seguro (CBC, Counter, OFB, CFB) sobre primitivas seguras.

%=====================================================================
\section{Un nuevo tipo de ataque: maleabilidad}
%=====================================================================

\begin{slidebox}
\textbf{Escenario.} Un sistema de RR.HH.\ guarda los sueldos de los empleados
\emph{cifrados} en una base de datos, con un criptosistema de flujo CPA-secure, para que ni
el administrador de la base pueda leerlos.
\end{slidebox}

\begin{center}
\renewcommand{\arraystretch}{1.3}
\begin{tabular}{c l}
\toprule
\textbf{empleado} & \textbf{sueldo (cifrado)}\\
\midrule
2542 & \texttt{0xEA26969AA3F61CDD9EC68498DF031CA935EFA9C}\\
2678 & \texttt{0xF4933E107178B88D8EE00F40E43A9A3C9D2EDF79}\\
2789 & \texttt{0x155BBBE7A5442CBBCAB8F57DACF7212255F3641C}\\
2890 & \texttt{0x4D9B47D518F2ED7DE04D601F1856767969A328E8}\\
\bottomrule
\end{tabular}
\end{center}

\textbf{La parte de cifrado es realmente segura.} El atacante no puede extraer información:
ni aproximar un sueldo, ni saber si es alto o bajo, ni sacar el número. Aunque le revelen que
\emph{uno} de los sueldos vale un millón, como el criptosistema es \textbf{no determinista},
puede cifrar el valor ``un millón'' un millón de veces y nunca obtendrá ninguno de los cuatro
textos cifrados de la tabla.

\subsection{El ataque: no hace falta conocer el valor}
A veces, \emph{para montar un ataque no se necesita conocer el valor}. Supóngase que el
atacante conoce su propio legajo (2678) y el de su jefe (2890), y que normalmente el jefe
cobra más. Entonces puede tomar el texto cifrado del jefe y hacer un \texttt{UPDATE}
copiándolo sobre su propia fila:

\begin{center}
\renewcommand{\arraystretch}{1.3}
\begin{tabular}{c l c}
\toprule
\textbf{empleado} & \textbf{sueldo (cifrado)} & \\
\midrule
2678 & \textcolor{attackRed}{\texttt{0x4D9B47D518F2ED7DE04D601F1856767969A328E8}}
     & \textcolor{attackRed}{\textbf{¡Aumento!}}\\
2890 & \texttt{0x4D9B47D518F2ED7DE04D601F1856767969A328E8} &\\
\bottomrule
\end{tabular}
\end{center}

\noindent\textbf{Aumento de sueldo instantáneo}, sin haber roto la confidencialidad. Una
solución ingenieril es cifrar el par \texttt{empleado}$\concat$\texttt{sueldo} y verificar, al
descifrar, que el legajo coincide con la fila. \emph{Pero esto no alcanza.}

\subsection{Modificaciones quirúrgicas: maleabilidad}

\begin{defbox}{Maleabilidad}
Un criptosistema es \textbf{maleable} cuando una modificación del \emph{texto cifrado} se
propaga de forma \textbf{sistemática y controlable} al \emph{texto plano}. En su forma más
simple: cambio un bit del cifrado y eso cambia un bit del plano.
\end{defbox}

Los cifrados de flujo son maleables. Si $\Enc_k(m)=G(k)\xor m=c$, y definimos $c'=c\xor x$,
entonces:
\[
  \Dec_k(c') \;=\; G(k)\xor c' \;=\; G(k)\xor(c\xor x)\;=\;\bigl(G(k)\xor c\bigr)\xor x \;=\; m\xor x.
\]
La modificación $x$ se traslada \emph{directamente} al plano. El atacante conoce el formato:

\begin{slidebox}
\textbf{Formato del campo:} $\;\mathrm{IV}\concat \mathrm{empl}\concat \mathrm{sueldo}$,
con $\mathrm{IV}=12$ bytes, $\mathrm{empl}=4$ bytes (int), $\mathrm{sueldo}=4$ bytes (int).
\end{slidebox}

\paragraph{Resolución del ejercicio del backstage.} El atacante conoce su propio sueldo en
plano (le llega el recibo): $1\,2345$. Quiere cambiarlo a $100\,000$ operando \emph{solo}
sobre el cifrado. Trabaja sobre los últimos 3 bytes (\texttt{2EDF79}). El procedimiento
iterativo (que puede colapsarse en un único XOR triple) es:

\begin{center}
\renewcommand{\arraystretch}{1.25}
\begin{tabular}{l l l l l}
\toprule
 & \multicolumn{3}{c}{representación binaria} & efecto al descifrar\\
\midrule
\texttt{2EDF79} & \texttt{00101110} & \texttt{11011111} & \texttt{01111001} & $\Dec_k(c)=12345$\\
$\xor\ 12345$  & \texttt{00000000} & \texttt{00110000} & \texttt{00111001} & $\Rightarrow \Dec_k(c')=0\ldots00$\\
$\xor\ 100000$ & \texttt{00000001} & \texttt{10000110} & \texttt{10100000} & $\Rightarrow \Dec_k(c'')=100000$\\
\midrule
$=$ \texttt{2F69F0} & \texttt{00101111} & \texttt{01101001} & \texttt{11110000} & (nuevo cifrado)\\
\bottomrule
\end{tabular}
\end{center}

\begin{itemize}[leftmargin=1.4em,itemsep=1pt]
  \item Primer paso ($\xor\,12345$): \emph{anula} el plano dejándolo en ceros
        ($12345\xor12345=0$). Es solo una ``plantilla'' de partida.
  \item Segundo paso ($\xor\,100000$): a partir de los ceros, fija el valor deseado
        ($0\xor100000=100000$).
\end{itemize}
El cifrado resultante \texttt{2F69F0} es una jeringoza que nadie entiende (no se tiene la
clave), pero al descifrarse da \textbf{exactamente} $100\,000$, y \emph{pasa toda la
verificación} del legajo. La maleabilidad es un problema de \textbf{integridad}, no
resuelto por el cifrado.

%=====================================================================
\section{Ataques de texto cifrado escogido (CCA)}
%=====================================================================

El formalismo que captura ``protegernos de cualquier modificación del cifrado'' es el
\textbf{ataque de texto cifrado escogido} (\emph{Chosen Ciphertext Attack}). Es una versión
más paranoica de CPA: el adversario tiene ahora \textbf{también un oráculo de descifrado}.

\begin{defbox}{Experimento $\mathrm{CCA}_{A,\Pi}$}
Dado un nivel de seguridad $n$, un adversario $A$ y un criptosistema $\Pi(n)$:
\begin{enumerate}[leftmargin=1.6em,itemsep=1pt]
  \item Se genera una clave $k\leftarrow \Kset$.
  \item $A$ obtiene $f(x)=e_k(x)$ \textbf{y} $g(x)=d_k(x)$, y genera $(m_0,m_1)$.
  \item Se genera $b\leftarrow\{0,1\}$.
  \item $A$ recibe $c=e_k(m_b)$, y \textbf{no puede calcular} $g(c)$
        (sino el ataque sería trivial).
  \item $A$ emite $b'\in\{0,1\}$.
\end{enumerate}
$\mathrm{CCA}_{A,\Pi}=1$ si $b=b'$. Si
$\Pr[\mathrm{CCA}_{A,\Pi}=1]<\tfrac12+\negl(n)$, $\Pi$ es \textbf{CCA-secure}.
\end{defbox}

% ---- Diagrama: oráculos CPA vs CCA ----
\begin{figure}[h]
\centering
\begin{tikzpicture}[font=\small,>=Stealth,node distance=5mm]
  \node[circle,draw,fill=itbaBlue,minimum size=1.1cm] (A) {$A$};
  \node[draw,rounded corners,fill=boxBlue,right=3.4cm of A,minimum width=2.6cm,
        minimum height=0.8cm] (enc) {$f(x)=e_k(x)$};
  \node[draw,rounded corners,fill=attackRed!20,below=0.7cm of enc,minimum width=2.6cm,
        minimum height=0.8cm] (dec) {$g(x)=d_k(x)$};
  \draw[-{Stealth}] (A.east) to[bend left=12] node[above,font=\scriptsize]{$x$} (enc.west);
  \draw[-{Stealth}] (enc.west) to[bend left=12] node[below,font=\scriptsize]{$e_k(x)$} (A.east);
  \draw[-{Stealth}] (A.east) to[bend right=12] node[above,font=\scriptsize,pos=0.6]{$x$} (dec.west);
  \draw[-{Stealth}] (dec.west) to[bend right=12] node[below,font=\scriptsize]{$d_k(x)$} (A.east);
  \node[font=\scriptsize,practGreen] at ($(enc.east)+(1.7,0)$) {oráculo de cifrado (CPA)};
  \node[font=\scriptsize,attackRed] at ($(dec.east)+(1.95,0)$) {oráculo de descifrado (solo CCA)};
\end{tikzpicture}
\caption{En CCA el adversario suma el \emph{oráculo de descifrado} $g$ al de cifrado $f$ del
escenario CPA. \textit{(Recreación conceptual de las diapositivas.)}}
\end{figure}

\begin{warnbox}{Ningún criptosistema visto es CCA-secure}
La buena noticia: CCA es una prueba \textbf{abarcativa} que cubre todas las familias de
modificaciones del cifrado (no hay que jugar al gato y al ratón). La mala: \textbf{ninguno}
de los criptosistemas vistos (ni los no vistos) la pasa, porque el cifrado está pensado para
\emph{confidencialidad} (transferencia de información), no para \emph{integridad}
(manipulación de la información).
\end{warnbox}

\begin{practbox}{Ejercicio (slide): el cifrado de flujo no es CCA-secure}
\textbf{Enunciado.} Demostrar que el cifrado de flujo en general no es CCA-secure.
\emph{Ayuda:} considerar $m_0=(0\ldots0)$ y $m_1=(1\ldots1)$ y verificar qué consultas se
pueden hacer a $g(x)$ al recibir $c$.

\medskip
\textbf{Resolución.} El adversario propone $m_0=0^n$ y $m_1=1^n$ y recibe el challenge
$c=e_k(m_b)=G(k)\xor m_b$. No puede preguntar $g(c)$, pero sí $g(c')$ para cualquier
$c'\neq c$. Construye $c'=c\xor 1^n$ (flip de todos los bits). Por maleabilidad:
\[
  g(c')=\Dec_k(c\xor 1^n)=m_b\xor 1^n=\overline{m_b}.
\]
Si $b=0$ entonces $m_b=0^n$ y $g(c')=1^n$; si $b=1$ entonces $g(c')=0^n$. El adversario lee
$g(c')$ y deduce $b$ con \textbf{probabilidad 1}. Como $c'\neq c$, la consulta es legal.
$\blacksquare$
\end{practbox}

Esto requiere \textbf{control de integridad}: identificar adulteraciones. La primitiva
criptográfica para ello es el \textbf{MAC} (\emph{Message Authentication Code}).

%=====================================================================
\section{Message Authentication Code (MAC)}
%=====================================================================

\begin{defbox}{Definición formal de MAC}
Un MAC es una \textbf{terna de algoritmos} $\langle\Gen,\Mac,\Vrfy\rangle$:
\[
  \Gen:()\to\Kset \quad(\text{generador de clave}),\qquad
  \Mac:\Kset\times\Pset\to\Tset \quad(\text{etiquetador}),
\]
\[
  \Vrfy:\Kset\times\Pset\times\Tset\to\{0,1\} \quad(\text{verificador}).
\]
\textbf{Propiedad básica:} para todo $m$ y $k$ válidos,
$\;\Vrfy_k\bigl(m,\Mac_k(m)\bigr)=1$.\\
Conjuntos: $\Kset$ (claves), $\Pset$ (planos), $\Tset$ (etiquetas/\emph{tags}).
\end{defbox}

El emisor envía el par $\langle m,t\rangle$ con $t=\Mac_k(m)$; el receptor calcula
$b=\Vrfy_k(m,t)$ y acepta si $b=1$. La función $\Mac$ parte de un mensaje plano arbitrario y
una clave y genera una etiqueta; $\Vrfy$ recibe clave, mensaje y etiqueta y devuelve un
booleano.

% ---- Diagrama MAC gen/verif (recreación slide práctica) ----
\begin{figure}[h]
\centering
\begin{tikzpicture}[font=\small,>=Stealth,node distance=6mm]
  % Emisor
  \node[font=\bfseries,itbaCyanDark] (em) at (0,2.4) {Emisor};
  \node[draw,fill=slideGray,minimum width=1.1cm,minimum height=0.7cm] (k1) at (-1.0,1.3) {$k$};
  \node[draw,fill=slideGray,minimum width=2.4cm,minimum height=0.7cm] (m1) at (1.2,1.3) {$m$};
  \node[draw,rounded corners,fill=itbaBlue,minimum width=1.4cm,minimum height=0.8cm] (mac) at (0.6,0.1) {$\Mac_k$};
  \node[draw,fill=practGreenL,minimum width=1.1cm,minimum height=0.7cm] (t1) at (0.6,-1.0) {$t$};
  \draw[-{Stealth}] (k1) |- (mac.west);
  \draw[-{Stealth}] (m1) -- (mac);
  \draw[-{Stealth}] (mac) -- (t1);
  % canal
  \node[draw,fill=practGreenL,minimum width=2.6cm,minimum height=0.8cm] (pair) at (5.2,0.1) {$\langle m,\,t\rangle$};
  \draw[-{Stealth},thick,practGreen] (t1.east) to[bend right=15] node[below,font=\scriptsize]{canal} (pair.south west);
  \draw[-{Stealth},thick,practGreen] (m1.east) to[bend left=10] (pair.north west);
  % Receptor
  \node[font=\bfseries,itbaCyanDark] (re) at (9.4,2.4) {Receptor};
  \node[draw,fill=slideGray,minimum width=1.1cm,minimum height=0.7cm] (k2) at (8.4,1.3) {$k$};
  \node[draw,rounded corners,fill=itbaBlue,minimum width=1.5cm,minimum height=0.8cm] (vrfy) at (9.4,0.1) {$\Vrfy_k$};
  \node[draw,fill=slideGray,minimum width=2.1cm,minimum height=0.7cm] (b) at (9.4,-1.0) {$b\in\{0,1\}$};
  \draw[-{Stealth}] (k2) -- (vrfy);
  \draw[-{Stealth}] (pair.east) to[bend left=12] (vrfy.west);
  \draw[-{Stealth}] (vrfy) -- (b);
\end{tikzpicture}
\caption{Generación y verificación de un MAC. El emisor adjunta $t=\Mac_k(m)$; el receptor
acepta si $\Vrfy_k(m,t)=1$. \textit{(Recreación de las diapositivas teórica y práctica.)}}
\end{figure}

\begin{warnbox}{Aclaraciones del profesor}
\begin{itemize}[leftmargin=1.3em,itemsep=2pt]
  \item \textbf{Un MAC \emph{no} provee no-repudio.} La clave $k$ es \textbf{simétrica}
        (compartida por emisor y receptor): cualquiera de los dos pudo generar la etiqueta,
        así que ante un tercero no se puede probar \emph{quién} la emitió. El no-repudio
        requiere firma digital (clave pública), que se ve más adelante.
  \item \textbf{Un MAC \emph{no} protege contra ataques de \emph{replay}.} Una etiqueta
        válida sigue siendo válida si se reenvía. Se mitiga incorporando un
        \textbf{número de secuencia} o un \textbf{timestamp} dentro del mensaje autenticado.
  \item \textbf{MAC ${\neq}$ dirección MAC} de redes: aquí es \emph{Message Authentication
        Code}.
\end{itemize}
\end{warnbox}

\subsection{Seguridad de un MAC: el experimento Mac-forge}

\begin{defbox}{Experimento $\mathrm{Mac\text{-}forge}_{A,\Pi}$}
Dado un nivel de seguridad $n$, un adversario $A$ y un esquema $\Pi(n)$:
\begin{enumerate}[leftmargin=1.6em,itemsep=1pt]
  \item Se genera una clave $k\leftarrow\Gen(1^n)$.
  \item $A$ recibe $1^n$ y acceso al \textbf{oráculo} $f(x)=\Mac_k(x)$.
        Realiza todas las evaluaciones que quiera (sea $Q$ el conjunto de consultas).
  \item $A$ emite un par $\langle m,t\rangle$ con $m\notin Q$.
\end{enumerate}
$\mathrm{Mac\text{-}forge}_{A,\Pi}=1$ (\textbf{éxito}) si y solo si
$\Vrfy_k(m,t)=1$ \textbf{y} $m\notin Q$.\\[2pt]
El MAC es \textbf{seguro (infalsificable)} ante un ataque de mensaje elegido si y solo si,
para todo adversario PPT $A$:
\[
  \Pr[\mathrm{Mac\text{-}forge}_{A,\Pi}(n)=1]\;\le\;\negl(n).
\]
\end{defbox}

% ---- Diagrama del experimento Mac-forge ----
\begin{figure}[h]
\centering
\begin{tikzpicture}[font=\small,>=Stealth]
  \node[circle,draw,fill=itbaBlue,minimum size=1.0cm] (A) at (0,0) {$A$};
  \node[draw,rounded corners,fill=boxBlue,minimum width=2.6cm,minimum height=0.9cm]
       (orac) at (4.2,0.6) {oráculo $\Mac_k(\cdot)$};
  \node[draw,rounded corners,fill=practGreenL,minimum width=2.6cm,minimum height=0.9cm]
       (out) at (4.2,-1.0) {emite $\langle m,t\rangle$};
  \draw[-{Stealth}] (A.east) to[bend left=10] node[above,font=\scriptsize]{$x_i\in Q$} (orac.west);
  \draw[-{Stealth}] (orac.west) to[bend left=10] node[below,font=\scriptsize]{$t_i$} (A.east);
  \draw[-{Stealth}] (A.south east) -- (out.west);
  \node[align=left,font=\scriptsize,anchor=west] at (7.0,-1.0)
       {gana si $\Vrfy_k(m,t)=1$\\ \textbf{y} $m\notin Q$};
\end{tikzpicture}
\caption{El adversario consulta libremente el oráculo (conjunto $Q$) y debe forjar una
etiqueta válida para un mensaje \textbf{nuevo} $m\notin Q$.}
\end{figure}

\begin{slidebox}
\textbf{Observaciones.} (i) El adversario puede obtener un MAC para \emph{cualquier} mensaje
que elija. (ii) Se considera \textbf{roto} el MAC si el adversario puede falsificar
cualquier mensaje, \textbf{independientemente de si tiene sentido o no}.
\end{slidebox}

%=====================================================================
\section{Ejercicio: ¿son seguros estos MACs?}
%=====================================================================

La siguiente serie de MACs ``ingenuos'' permite extraer las \textbf{propiedades que un MAC
seguro debe cumplir}. Todos se rompen con Mac-forge. Resolvemos cada uno con el paso a paso
de la clase.

\begin{practbox}{(a) $\Mac_k(m)=G(k)\xor m$ — usar un cifrado de flujo como MAC}
$G$ es un generador pseudoaleatorio: la etiqueta es tan larga como el mensaje y resulta de
hacer XOR contra la secuencia $G(k)$.

\medskip
\textbf{Problema:} la salida $G(k)$ es \emph{determinística} (con la misma clave, el primer
bloque del generador es siempre el mismo). Etiquetar es simplemente hacer XOR contra $G(k)$,
así que si el adversario \textbf{averigua $G(k)$ puede etiquetar cualquier mensaje}.

\medskip
\textbf{Algoritmo atacante (Mac-forge):}
\begin{enumerate}[leftmargin=1.6em,itemsep=1pt]
  \item Consulta el oráculo con $m=0\ldots0$. Obtiene $t_0=\Mac_k(0)=G(k)\xor 0=\boxed{G(k)}$.
        \emph{¡Recuperó la corriente de clave!}
  \item Elige un mensaje nuevo $m'\neq 0$, por ejemplo $m'=1\ldots1$, y calcula
        $t'=G(k)\xor m'$.
  \item Emite $\langle m', t'\rangle$ con $m'\notin Q$: pasa la verificación. $\blacksquare$
\end{enumerate}
\end{practbox}

\begin{practbox}{(b) $\Mac_k(m)=k\xor \mathrm{first\_k\_bits}(m)$}
La etiqueta es el XOR de la clave (longitud fija) con los \emph{primeros} $|k|$ bits del
mensaje. \emph{Hay dos falsificaciones distintas:}

\medskip
\textbf{Ataque 1 — recuperar la clave (la del curso).}
\begin{enumerate}[leftmargin=1.6em,itemsep=1pt]
  \item Consulta $m=0\ldots0$: $\;t=k\xor \mathrm{first\_k\_bits}(0)=k\xor 0=\boxed{k}$.
        \emph{¡Recuperó la clave entera!}
  \item Emite un mensaje nuevo, p.\ ej.\ $m'=1\ldots1$ con etiqueta
        $t'=k\xor(\underbrace{1\ldots1}_{k\text{ bits}})$. Pasa la verificación. $\blacksquare$
\end{enumerate}

\medskip
\textbf{Ataque 2 — más fácil aún (la lección buscada).} El MAC \textbf{ignora} todo el
mensaje salvo sus primeros $|k|$ bits. Entonces:
\begin{enumerate}[leftmargin=1.6em,itemsep=1pt]
  \item Consulta un mensaje $m$ y obtiene su etiqueta $t$.
  \item Emite $\langle m\concat\texttt{<carácter extra>},\ t\rangle$: el sufijo agregado cae
        fuera de los primeros $|k|$ bits, así que la \textbf{misma etiqueta} sigue siendo
        válida para un mensaje distinto. $\blacksquare$
\end{enumerate}
\textbf{Lección:} \emph{un MAC que no considera la totalidad del mensaje es trivialmente
falsificable.} Si ignora algún bit, basta modificar ese bit y reusar la etiqueta.
\end{practbox}

\begin{practbox}{(c) $\Mac_k(m)=\Enc_k(|m|)$ con $\Enc$ CPA-secure}
La etiqueta cifra \textbf{solo la longitud} del mensaje, no su contenido.

\medskip
\textbf{Verificación:} hay que descifrar la etiqueta y ver si la longitud coincide.
\textbf{Problema:} la etiqueta \emph{no depende del contenido}.

\medskip
\textbf{Algoritmo atacante:}
\begin{enumerate}[leftmargin=1.6em,itemsep=1pt]
  \item Pide el etiquetado de un mensaje $m_A$ de longitud $L$: obtiene $t_A$.
  \item Toma cualquier \emph{otro} mensaje $m_B\neq m_A$ de la \textbf{misma longitud} $L$.
  \item Emite $\langle m_B, t_A\rangle$. Al verificar, se descifra $t_A\Rightarrow L$, que
        coincide con $|m_B|$: \textbf{válido}. $\blacksquare$
\end{enumerate}
\emph{(Como $\Enc$ es no determinista, $t_A\neq t_B$ aunque $|m_A|=|m_B|$; pero ambos
descifran a la misma longitud, así que las etiquetas son intercambiables.)}
\end{practbox}

\begin{slidebox}
\textbf{Conclusión.} Un MAC robusto/infalsificable es \emph{siempre} el resultado de
procesar la \textbf{totalidad del contenido} del mensaje, sin determinismos triviales
extraíbles ni partes ignoradas.
\end{slidebox}

%=====================================================================
\section{Cómo construir un MAC: CBC-MAC}
%=====================================================================

Primera familia de construcciones: \textbf{a partir de una primitiva de cifrado de bloque}.
Sea $F$ una función pseudoaleatoria ($F_k:\{0,1\}^n\to\{0,1\}^n$).

\subsection{Construcción para mensajes de longitud fija}

\begin{defbox}{MAC seguro para mensajes de longitud fija $n$}
\textbf{Gen:} $k\leftarrow\{0,1\}^n$. \quad
\textbf{Mac:} $t\leftarrow F_k(m)$ con $|m|=|k|=n$. \quad
\textbf{Vrfy:} si $|m|\neq|k|\Rightarrow 0$; si $F_k(m)==t\Rightarrow 1$.
\end{defbox}

Como $F_k$ es una PRF de $n$ bits a $n$ bits que procesa \emph{todo} el bloque, esta
construcción es infalsificable \textbf{para mensajes de longitud fija}.

\subsection{CBC-MAC (para mensajes más largos, longitud fija)}

\begin{defbox}{CBC-MAC}
\textbf{Gen:} $k\leftarrow\{0,1\}^n$.\quad \textbf{Mac:} dado $m=m_1\concat m_2\concat
\cdots\concat m_\ell$ (bloques de $n$ bits), con $t_0=0^n$:
\[
  t_i = F_k\bigl(t_{i-1}\xor m_i\bigr),\qquad \Mac_k(m)=t_\ell.
\]
\textbf{Vrfy:} $\Vrfy_k(m,t)=1 \iff t=\Mac_k(m)$.
\end{defbox}

A diferencia del CBC \emph{de cifrado}, aquí se emite \textbf{una sola etiqueta} (la del
último bloque) y \textbf{no hay IV} (el estado inicial es fijo en $0^n$). Se encadena un
bloque con el siguiente; importa el orden y los estados intermedios $t_1,\dots,t_{\ell-1}$
\textbf{se descartan}.

% ---- Diagrama CBC-MAC ----
\begin{figure}[h]
\centering
\begin{tikzpicture}[font=\small,>=Stealth,node distance=4mm,
   blk/.style={draw,fill=slideGray,minimum width=1.2cm,minimum height=0.7cm},
   F/.style={draw,rounded corners,fill=itbaBlue,minimum width=1.0cm,minimum height=0.8cm},
   xo/.style={circle,draw,inner sep=1.2pt}]
  % mensajes
  \node[blk] (m1) at (0,2.4) {$m_1$};
  \node[blk] (m2) at (2.4,2.4) {$m_2$};
  \node[blk] (md) at (4.8,2.4) {$\cdots$};
  \node[blk] (ml) at (7.2,2.4) {$m_\ell$};
  % t0
  \node[blk] (t0) at (-2.0,0.9) {$0^n$};
  % xor row
  \node[xo] (x1) at (0,0.9) {$\xor$};
  \node[xo] (x2) at (2.4,0.9) {$\xor$};
  \node[xo] (x3) at (4.8,0.9) {$\xor$};
  \node[xo] (x4) at (7.2,0.9) {$\xor$};
  % F row
  \node[F] (F1) at (0,-0.4) {$F_k$};
  \node[F] (F2) at (2.4,-0.4) {$F_k$};
  \node[draw,minimum width=1.0cm,minimum height=0.8cm] (F3) at (4.8,-0.4) {$\cdots$};
  \node[F] (F4) at (7.2,-0.4) {$F_k$};
  % outputs
  \node[font=\scriptsize] (o0) at (-2.0,-1.4) {$t_0$};
  \node[font=\scriptsize] (o1) at (0,-1.4) {$t_1$};
  \node[font=\scriptsize] (o2) at (2.4,-1.4) {$t_2$};
  \node[draw,fill=practGreenL,minimum width=1.0cm] (tl) at (7.2,-1.6) {$t_\ell$};
  % wires
  \draw[-{Stealth}] (m1)--(x1); \draw[-{Stealth}] (m2)--(x2);
  \draw[-{Stealth}] (md)--(x3); \draw[-{Stealth}] (ml)--(x4);
  \draw[-{Stealth}] (t0)--(x1);
  \draw[-{Stealth}] (x1)--(F1); \draw[-{Stealth}] (x2)--(F2);
  \draw[-{Stealth}] (x3)--(F3); \draw[-{Stealth}] (x4)--(F4);
  \draw[-{Stealth}] (F1) |- (o1.west) ; \draw[-{Stealth}] (F1.east) -| (x2);
  \draw[-{Stealth}] (F2.east) -| (x3);
  \draw[-{Stealth}] (F3.east) -| (x4);
  \draw[-{Stealth}] (F4) -- (tl);
\end{tikzpicture}
\caption{CBC-MAC: encadenamiento $t_i=F_k(t_{i-1}\xor m_i)$ con $t_0=0^n$; la etiqueta es el
\textbf{último} $t_\ell$. \textit{(Recreación de las diapositivas.)}}
\end{figure}

\begin{warnbox}{Aclaración clave del profesor}
La construcción CBC-MAC anterior es \textbf{infalsificable SOLO si se permiten mensajes de
una misma longitud}. Con mensajes de \textbf{longitud variable} es falsificable.
\end{warnbox}

\subsection{Ataque a CBC-MAC con longitud variable}
El atacante explota que los \textbf{estados intermedios} ($t_1,\dots$) se descartan, pero
puede recalcularlos: el estado tras procesar dos bloques, $F_k(\cdot)\xor m_3$, le permite
``empalmar'' mensajes.

\begin{practbox}{Ataque (slide teórica): forja con longitud variable}
\begin{enumerate}[leftmargin=1.6em,itemsep=2pt]
  \item Crear dos bloques aleatorios $A$ y $B$. Definir
        $m_1=A\concat B$ \;y\; $m_2=A$.
  \item Consultar el oráculo: $t_1=f(m_1)$, $\;t_2=f(m_2)$.
  \item \textbf{Observar:} $t_2=F_k(0^n\xor A)=F_k(A)$, y
        $t_1=F_k\bigl(F_k(A)\xor B\bigr)=F_k(t_2\xor B)$.
  \item Emitir el mensaje de \textbf{tres} bloques
        $\;m^\ast = A\concat B\concat (A\xor t_1)\;$ con etiqueta $\;t_2$.
\end{enumerate}
\textbf{Por qué funciona.} Calculamos $\Mac_k(m^\ast)$ paso a paso:
\[
\begin{aligned}
  s_1 &= F_k(0^n\xor A)=F_k(A)=t_2,\\
  s_2 &= F_k(s_1\xor B)=F_k(t_2\xor B)=t_1,\\
  s_3 &= F_k\bigl(s_2\xor (A\xor t_1)\bigr)=F_k\bigl(t_1\xor A\xor t_1\bigr)=F_k(A)=t_2.
\end{aligned}
\]
Luego $\Mac_k(m^\ast)=s_3=t_2$, que es justamente la etiqueta emitida, y $m^\ast\notin Q$.
\textbf{Falsificación válida.} $\blacksquare$
\end{practbox}

\subsection{Extensiones seguras de CBC-MAC para longitud arbitraria}
Para autenticar mensajes de \textbf{cualquier} longitud existen tres opciones seguras (y una
insegura para descartar):

\begin{defbox}{Opciones seguras}
\begin{description}[leftmargin=2.2em,style=nextline,itemsep=4pt]
  \item[Opción 1 — derivar clave de la longitud]
        $\;k'=F_k(|m|)\;$ y $\;t\leftarrow \text{CBC-MAC}_{k'}(m).$
        \emph{(Cada longitud usa una clave distinta.)}
  \item[Opción 2 — longitud como \underline{prefijo}]
        $\;m'=|m|\concat m\;$ y $\;t\leftarrow\text{CBC-MAC}_k(m').$
        \emph{La longitud debe ir al principio, no al final.}
  \item[Opción 3 — doble clave]
        $\;k_1,k_2\leftarrow\{0,1\}^n;\;\;t\leftarrow\text{CBC-MAC}_{k_1}(m);\;\;
        \hat t\leftarrow F_{k_2}(t).$
\end{description}
\end{defbox}

\begin{warnbox}{La longitud como \emph{sufijo} NO es segura}
Si $m'=m\concat |m|$ y $\Mac_k(m)=\text{CBC-MAC}_k(m')$, hay un ataque. Con bloques
aleatorios $A,B,C$:
\[
\begin{aligned}
  m_1=AAA &\to m_1'=AAA3, & t_1&=f(m_1)\\
  m_2=BBB &\to m_2'=BBB3, & t_2&=f(m_2)\\
  m_3=AAA3CC &\to m_3'=AAA3CC6, & t_3&=f(m_3)\\
  &\text{Sea } X=t_1\xor t_2\xor C; &&\text{emitir } (BBB3XC,\ t_3).
\end{aligned}
\]
La etiqueta $t_3$ resulta válida para un mensaje no consultado. \emph{Por eso la longitud va
como prefijo.} \textit{(Notación de la slide; $3,6$ representan los sufijos de longitud.)}
\end{warnbox}

%=====================================================================
\section{Práctica: construcciones para mensajes de longitud variable}
%=====================================================================

La clase práctica desarrolla \textbf{en detalle} por qué las construcciones ``naturales''
para extender un MAC de bloque fijo $\Mac'$ a mensajes de longitud variable fallan, antes de
llegar a la solución (ineficiente pero segura) y, finalmente, a CBC-MAC.

\begin{practbox}{Sugerencia 1 — XOR de los bloques y autenticar el resultado}
\[
  t := \Mac'_k\Bigl(\bigoplus_i m_i\Bigr),\qquad \text{emitir } \langle m_1,m_2,\ldots,m_d;\ t\rangle.
\]
\textbf{Falla.} El adversario puede falsificar un tag válido sobre un \emph{nuevo} mensaje
cambiando el mensaje original \textbf{de modo que el XOR de los bloques no cambie}. Por
ejemplo, mueve un bit de un bloque a otro: si el XOR total da lo mismo, la etiqueta sigue
siendo válida y $m'\neq m$. $\blacksquare$
\end{practbox}

\begin{practbox}{Sugerencia 2 — autenticar cada bloque por separado}
\[
  t_i := \Mac'_k(m_i),\qquad \text{emitir } \langle m_1,m_2,\ldots,m_d;\ t_1,t_2,\ldots,t_d\rangle.
\]
\textbf{Falla.} El adversario consulta cada bloque por separado ($\Mac'$ de $m_1$, de
$m_2$\ldots) y luego \textbf{reordena/permuta} los bloques con sus etiquetas. Por ejemplo, de
$\langle m_1,m_2,\ldots; t_1,t_2,\ldots\rangle$ arma
$\langle m_2,m_1,m_3; t_2,t_1,t_3\rangle$: cada par $(m_i,t_i)$ valida bien, así que el
mensaje permutado pasa. También puede \textbf{mezclar bloques de mensajes distintos}.
$\blacksquare$
\end{practbox}

\begin{practbox}{Sugerencia 3 — autenticar cada bloque con su número de secuencia}
\[
  t_i := \Mac'_k(i\concat m_i),\qquad \text{emitir } \langle m_1,\ldots,m_d;\ t_1,\ldots,t_d\rangle.
\]
\textbf{Mejora pero aún falla.} El número de secuencia impide reordenar (ahora cada $t_i$
depende de la posición $i$). \textbf{Pero} el adversario todavía puede \textbf{mezclar
bloques de mensajes diferentes que estén en la misma posición}: si consulta dos mensajes y
toma $m_1$ del primero, $m'_2$ del segundo, $m_3$ del primero\ldots\ cada $(i\concat m)$
sigue correspondiendo a su tag válido $\Rightarrow$ el mensaje híbrido valida. $\blacksquare$
\end{practbox}

\begin{practbox}{Construcción segura (ineficiente): random + longitud + secuencia + bloque}
A partir de $\Mac'=\langle\Gen',\Mac',\Vrfy'\rangle$ para mensajes de longitud $n$:
\begin{itemize}[leftmargin=1.4em,itemsep=1pt]
  \item \textbf{Gen:} $\Gen'$, con $|k|=n$.
  \item \textbf{Mac:} para $|m|=L<2^{n/4}$, se parte $m$ en $d$ bloques de $n/4$ bits (se
        completa con ceros) y se elige $r\leftarrow\{0,1\}^{n/4}$:
        \[
          t_i = \Mac'_k\bigl(r\concat L\concat i\concat m_i\bigr),\quad i=1,\ldots,d,
          \qquad \text{emitir } \langle r,\,t_1,\ldots,t_d\rangle.
        \]
\end{itemize}
\textbf{Por qué es seguro:}
\begin{itemize}[leftmargin=1.4em,itemsep=1pt]
  \item \textbf{El número $r$ aleatorio} es el \emph{mismo} para todos los bloques de un mismo
        mensaje y \emph{distinto} entre mensajes. No se pueden \textbf{mezclar bloques de
        mensajes distintos}: tendrían valores de $r$ que no coinciden y la verificación falla
        (al armar un mensaje híbrido, ¿qué $r$ emitir en el tag? No hay uno que sirva para
        todos).
  \item \textbf{La longitud $L$} hace saber cuántos bloques deben llegar (impide
        truncamientos/extensiones).
  \item \textbf{El número de secuencia $i$} mantiene el orden (impide permutar).
\end{itemize}
\textbf{Problema: ¡ineficiente!} Emite \emph{todos} los tags por separado (una barbaridad de
pedacitos) y además parte el mensaje en bloques de $n/4$. Es seguro pero poco práctico
$\Rightarrow$ se usa \textbf{CBC-MAC}, que emite una sola etiqueta.
\end{practbox}

%=====================================================================
\section{Funciones de hash criptográficas}
%=====================================================================

\begin{defbox}{Función de hash}
Son \textbf{pares de algoritmos} $\langle\Gen,H\rangle$:
\[
  \Gen:s\leftarrow\Sset,\qquad h=H^s(m)\in\{0,1\}^L,
\]
con $L$ la longitud (\textbf{fija}) del hash. \textbf{Comprimen} una entrada de longitud
arbitraria a una salida de longitud fija (por eso, \emph{funciones de resumen / digest}).
\end{defbox}

\begin{slidebox}
\textbf{Diferencia importante con un MAC.} $s$ \textbf{no es una clave}: es un mero
\emph{selector} (público) de la función dentro de la familia. En muchas implementaciones
$\Sset=\{s_0\}$ (hay una única función). Las funciones de hash son \textbf{análogas a los
MACs, pero sin clave}.
\end{slidebox}

\subsection{Colisiones y la propiedad fundamental}
\begin{defbox}{Colisión}
Una \textbf{colisión} es un par $x\neq x'$ con $H(x)=H(x')$.\\
\textbf{Principio del palomar:} si hay $n+1$ mensajes y $n$ valores de salida posibles,
existe al menos una colisión. \emph{Las colisiones van a existir siempre} (la entrada es más
grande que la salida); el objetivo es que sean \textbf{computacionalmente difíciles} de
encontrar.
\end{defbox}

\subsection{Las tres propiedades (niveles de seguridad)}
\begin{center}
\renewcommand{\arraystretch}{1.3}
\begin{tabular}{>{\bfseries}p{4.2cm} p{9.2cm}}
\toprule
Resistencia a preimágenes & Dado $y$, es computacionalmente imposible hallar $x$ tal que
$h(x)=y$. (\emph{One-way}.)\\
Resistencia a segundas preimágenes & Dado $x$, es imposible hallar $x'\neq x$ con
$h(x')=h(x)$.\\
Resistencia a colisiones & Es imposible hallar \emph{cualquier} par $x,x'$ con $h(x)=h(x')$ y
$x\neq x'$.\\
\bottomrule
\end{tabular}
\end{center}

% ---- Diagrama: los tres niveles de seguridad ----
\begin{figure}[h]
\centering
\begin{tikzpicture}[font=\scriptsize,>=Stealth,
   dom/.style={draw,ellipse,minimum width=2.0cm,minimum height=2.6cm},
   img/.style={draw,ellipse,minimum width=1.6cm,minimum height=2.6cm}]
  % colisiones
  \begin{scope}[xshift=0cm]
    \node[font=\bfseries\normalsize,itbaCyanDark] at (0.6,2.0) {Colisiones};
    \node[dom] (d) at (0,0) {}; \node[img] (i) at (2.4,0) {};
    \fill (d.center)++(0,0.5) circle (1.6pt) coordinate (a);
    \fill (d.center)++(0,-0.5) circle (1.6pt) coordinate (b);
    \fill (i.center) circle (1.6pt) coordinate (c);
    \draw[-{Stealth}] (a)--(c); \draw[-{Stealth}] (b)--(c);
    \node at (0.6,-1.9) {hallar $x\neq x'$};
  \end{scope}
  % segundas preimagenes
  \begin{scope}[xshift=5.4cm]
    \node[font=\bfseries\normalsize,itbaCyanDark] at (0.9,2.0) {2.ª preimagen};
    \node[dom] (d2) at (0,0) {}; \node[img] (i2) at (2.4,0) {};
    \fill[attackRed] (d2.center)++(0,0.5) circle (1.8pt) coordinate (a2);
    \node[attackRed] at ($(d2.center)+(-0.55,0.5)$) {$x$ dado};
    \fill (d2.center)++(0,-0.5) circle (1.6pt) coordinate (b2);
    \fill (i2.center) circle (1.6pt) coordinate (c2);
    \draw[-{Stealth},attackRed] (a2)--(c2); \draw[-{Stealth}] (b2)--(c2);
    \node at (0.9,-1.9) {dado $x$, hallar $x'$};
  \end{scope}
  % preimagenes
  \begin{scope}[xshift=10.8cm]
    \node[font=\bfseries\normalsize,itbaCyanDark] at (0.9,2.0) {Preimagen};
    \node[dom] (d3) at (0,0) {}; \node[img] (i3) at (2.4,0) {};
    \fill (d3.center) circle (1.6pt) coordinate (a3);
    \fill[attackRed] (i3.center) circle (1.8pt) coordinate (c3);
    \node[attackRed] at ($(i3.center)+(0.0,0.9)$) {$y$ dado};
    \draw[-{Stealth}] (a3)--(c3);
    \node at (0.9,-1.9) {dado $y$, hallar $x$};
  \end{scope}
\end{tikzpicture}
\caption{Los tres niveles de seguridad. La \textbf{resistencia a colisiones} es donde el
atacante tiene \emph{mayor} ventaja (mayor grado de libertad: elige ambos $x,x'$); implica
las otras dos. \textit{(Recreación de las diapositivas.)}}
\end{figure}

\begin{defbox}{Experimento $\mathrm{Hash\text{-}Coll}_{A,H}$}
Dado un nivel $n$, un adversario $A$ y una familia $H^s(n)$:
\begin{enumerate}[leftmargin=1.6em,itemsep=1pt]
  \item Se selecciona una función $s\leftarrow\Sset$.
  \item $A$ obtiene acceso a $H(x)=H^s(x)$.
  \item $A$ emite $x,x'$.
\end{enumerate}
$\mathrm{Hash\text{-}Coll}_{A,H}=1$ si $x\neq x'$ y $H(x)=H(x')$. Si
$\Pr[\mathrm{Hash\text{-}Coll}_{A,H}=1]<\negl(n)$, $H$ es \textbf{libre de colisiones}.
\textbf{El desafío:} $s$ se fija en el momento, nada precalculado, y en tiempo polinómico hay
que materializar dos mensajes que colisionen.
\end{defbox}

\subsection{La paradoja del cumpleaños}
\begin{defbox}{Costo de los ataques por fuerza bruta. Sea $h:A\to B$.}
\begin{itemize}[leftmargin=1.4em,itemsep=2pt]
  \item \textbf{Preimágenes:} requiere $|B|$ intentos (en promedio).
  \item \textbf{Segundas imágenes:} requiere $|B|$ intentos.
  \item \textbf{Colisiones:} requiere solo $|B|^{1/2}$ intentos. \textbf{Paradoja del
        cumpleaños.}
\end{itemize}
\end{defbox}

\textbf{¿Por qué la raíz cuadrada?} Buscando una colisión arbitraria, el atacante genera un
mensaje, lo compara contra \emph{todos} los anteriores; genera otro y lo compara con todos
los previos\ldots\ Cada nuevo mensaje se chequea contra el conjunto acumulado, de modo que la
cantidad de \emph{pares} crece cuadráticamente. La complejidad cae a
$\sqrt{|B|}$: medido en bits, es \textbf{la mitad del nivel de seguridad} que las otras
propiedades. Es la misma optimización que en probabilidad: una cosa es la chance de que
alguien coincida con \emph{una} fecha fija, otra es que \emph{cualquiera} coincida con
\emph{cualquiera} en un grupo.

% ---- Diagrama paradoja del cumpleaños ----
\begin{figure}[h]
\centering
\begin{tikzpicture}[font=\scriptsize,>=Stealth]
  \foreach \i/\x in {1/0,2/1.2,3/2.4,4/3.6,5/4.8}{
     \node[draw,circle,fill=itbaBlue,minimum size=0.55cm] (n\i) at (\x,0) {$x_\i$};
  }
  \node[font=\itshape] at (6.3,0) {\dots};
  \foreach \a/\b in {1/2,1/3,2/3,1/4,2/4,3/4,1/5,2/5,3/5,4/5}{
     \draw[gray!55] (n\a) to[bend left=22] (n\b);
  }
  \node[align=center,font=\small] at (2.4,-1.6)
     {Cada nuevo mensaje se compara con \emph{todos} los anteriores:\\
      $\sim\binom{q}{2}$ pares $\Rightarrow$ colisión esperada con $q\approx\sqrt{|B|}$.};
\end{tikzpicture}
\caption{Paradoja del cumpleaños: la cantidad de pares crece como $q^2$, reduciendo el costo
de hallar una colisión a $\sqrt{|B|}$. \textit{(Recreación conceptual.)}}
\end{figure}

\subsection{El modelo iterativo: Merkle-Damgård}
Propuesto por \textbf{Merkle (1989)} y usado por muchas funciones. Procesa entradas de
longitud arbitraria a partir de una función de \textbf{compresión} $h^s$ de bloque pequeño.

\begin{itemize}[leftmargin=1.4em,itemsep=1pt]
  \item \textbf{Preprocesamiento:} ajustar el tamaño del mensaje (\emph{padding}) y agregar
        un bloque con la \textbf{longitud}.
  \item \textbf{Función de compresión $f$:} similar al hash pero sobre bloques pequeños;
        \textbf{la seguridad del esquema está dada por la función de compresión} (si $h^s$ es
        mala, todo el esquema es malo).
\end{itemize}

% ---- Diagrama Merkle-Damgard ----
\begin{figure}[h]
\centering
\begin{tikzpicture}[font=\scriptsize,>=Stealth,
   tr/.style={draw,trapezium,trapezium left angle=70,trapezium right angle=110,
              fill=itbaBlue!60,minimum width=1.6cm,minimum height=1.0cm,shape border rotate=180},
   blk/.style={draw,fill=slideGray,minimum width=1.7cm,minimum height=0.55cm},
   nout/.style={draw,fill=white,minimum width=0.9cm,minimum height=0.5cm}]
  \node[font=\itshape] (z0) at (-1.9,1.5) {$z_0=0^\ell$};
  % three compression stages
  \node[blk] (b1) at (0,1.5) {$0^\ell\concat x_1$};
  \node[tr]  (h1) at (0,0.4) {$h^s$};
  \node[nout] (o1) at (0,-0.9) {$z_1$};
  \node[blk] (b2) at (2.9,1.5) {$z_1\concat x_2$};
  \node[tr]  (h2) at (2.9,0.4) {$h^s$};
  \node[nout] (o2) at (2.9,-0.9) {$z_2$};
  \node at (5.0,0.4) {$\cdots$};
  \node[blk] (b3) at (7.6,1.5) {$z_{B}\concat x_{B+1}$};
  \node[tr]  (h3) at (7.6,0.4) {$h^s$};
  \node[nout,fill=practGreenL] (o3) at (7.6,-0.9) {$z_{B+1}$};
  \node[font=\itshape] at (7.6,2.3) {$x_{B+1}=|x|$};
  % wires
  \draw[-{Stealth}] (b1)--(h1); \draw[-{Stealth}] (h1)--(o1);
  \draw[-{Stealth}] (b2)--(h2); \draw[-{Stealth}] (h2)--(o2);
  \draw[-{Stealth}] (b3)--(h3); \draw[-{Stealth}] (h3)--(o3);
  \draw[-{Stealth}] (o1) -| (1.45,1.1) |- (b2.west);
  \draw[-{Stealth}] (o2) -| (4.3,1.1) -- (4.9,1.1);
  \draw[-{Stealth}] (5.3,1.1) -| (6.1,1.1) |- (b3.west);
  \node[draw,fill=practGreenL,minimum width=1.0cm] at (9.3,-0.9) {$H^s(x)$};
  \draw[-{Stealth}] (o3) -- (8.75,-0.9);
\end{tikzpicture}
\caption{Transformación de Merkle-Damg\aa rd: $z_i=h^s(z_{i-1}\concat x_i)$, con $z_0=0^\ell$
y el bloque final $x_{B+1}=|x|$. \textbf{Una colisión en $H^s$ solo puede ocurrir si hay una
colisión en $h^s$.} \textit{(Recreación de las diapositivas.)}}
\end{figure}

\subsection{Primitivas concretas}
\begin{center}
\renewcommand{\arraystretch}{1.25}
\begin{tabular}{>{\bfseries}l c p{6.6cm}}
\toprule
Primitiva & Salida & Notas \\
\midrule
MD5   & 128 bits & \textcolor{attackRed}{Quebrada}; aún aparece en protocolos viejos.\\
SHA-1 & 160 bits & Se construye sobre MD5; \emph{al borde}, complicada.\\
SHA-2 & 256/384/512 bits & Estándar durante mucho tiempo.\\
SHA-3 & 224/256/384/512 bits & Estandarizado por NIST (2013), antes \emph{Keccak}; usa
\textbf{modelo esponja}. \textbf{Recomendado} para proyectos nuevos.\\
\bottomrule
\end{tabular}
\end{center}

\begin{slidebox}
\textbf{Ejemplos (entrada $\to$ digest hex).}\\[2pt]
{\footnotesize\ttfamily
MD5:\ \ ""\,$\to$\,d41d8cd98f00b204e9800998ecf8427e;\quad "abc"\,$\to$\,900150983cd24fb0d6963f7d28e17f72\\
SHA-1: ""\,$\to$\,da39a3ee5e6b4b0d3255bfef95601890afd80709;\ "abc"\,$\to$\,a9993e364706816aba3e25717850c26c9cd0d89d}
\end{slidebox}

\begin{warnbox}{Tamaño mínimo en la práctica}
El tamaño mínimo de salida usado hoy es \textbf{160 bits}, para que la búsqueda de colisiones
requiera más de $2^{80}$ operaciones (preimágenes/2.ª imágenes: $2^{160}$).
\textbf{Condición necesaria pero NO suficiente:} aunque la salida sea grande, si la función
es mala el esquema es inseguro. Las funciones se miden por el nivel de seguridad de la
\emph{resistencia a colisiones} (la mitad de los bits).
\end{warnbox}

%=====================================================================
\section{MACs a partir de funciones de hash: HMAC y NMAC}
%=====================================================================

Segunda familia de construcciones de MAC: \textbf{a partir de una función de hash} libre de
colisiones. Ventaja sobre CBC-MAC: las funciones de hash están diseñadas para procesar
volúmenes grandes de datos muy rápido. Un MAC basado en hash suele ser \textbf{1 o 2 órdenes
de magnitud más rápido} que un CBC-MAC.

\subsection{NMAC (Nested MAC)}
\begin{defbox}{NMAC}
\textbf{Gen:} emite $(s,k_1,k_2)$ con $|k_1|=|k_2|=n$.\quad
\textbf{Mac:} para $m=m_1m_2\cdots m_B$ con $|m|=L$,
\[
  t := h^s_{k_1}\bigl(H^s_{k_2}(m)\bigr).
\]
\end{defbox}
Se hace el hash con clave del mensaje (clave $k_2$ y la longitud $L$ como bloque final)
obteniendo $z=H^s_{k_2}(m)$, y luego un \emph{último} paso con la otra clave $k_1$:
$t=h^s_{k_1}(z)$.

% ---- Diagrama NMAC ----
\begin{figure}[h]
\centering
\begin{tikzpicture}[font=\scriptsize,>=Stealth,
   tr/.style={draw,trapezium,trapezium left angle=70,trapezium right angle=110,
              fill=itbaBlue!50,minimum width=1.5cm,minimum height=0.95cm,shape border rotate=180},
   blk/.style={draw,fill=slideGray,minimum width=1.5cm,minimum height=0.5cm},
   nout/.style={draw,fill=white,minimum width=0.8cm,minimum height=0.45cm}]
  \node[blk](b1) at (0,1.4){$k_2\concat m_1$};
  \node[tr](h1) at (0,0.35){$h^s$}; \node[nout](z1) at (0,-0.75){$z_1$};
  \node[blk](b2) at (2.6,1.4){$z_1\concat m_2$};
  \node[tr](h2) at (2.6,0.35){$h^s$}; \node[nout](z2) at (2.6,-0.75){$z_2$};
  \node at (4.3,0.35){$\cdots$};
  \node[blk](b3) at (6.0,1.4){$z_{B-1}\concat m_B$};
  \node[tr](h3) at (6.0,0.35){$h^s$}; \node[nout](zB) at (6.0,-0.75){$z_B$};
  \node[blk](b4) at (8.4,1.4){$z_B\concat L$};
  \node[tr](h4) at (8.4,0.35){$h^s$}; \node[nout,fill=itbaBlue!25](z) at (8.4,-0.75){$z$};
  \node[blk](b5) at (10.8,1.4){$k_1\concat z$};
  \node[tr](h5) at (10.8,0.35){$h^s$};
  \node[draw,fill=practGreenL,minimum width=0.8cm](t) at (10.8,-0.75){$t$};
  \foreach \b/\h in {b1/h1,b2/h2,b3/h3,b4/h4,b5/h5}{\draw[-{Stealth}](\b)--(\h);}
  \draw[-{Stealth}](h1)--(z1);\draw[-{Stealth}](h2)--(z2);\draw[-{Stealth}](h3)--(zB);
  \draw[-{Stealth}](h4)--(z);\draw[-{Stealth}](h5)--(t);
  \draw[-{Stealth}](z1)-|(1.3,1.05)|-(b2.west);
  \draw[-{Stealth}](zB)-|(7.2,1.05)|-(b4.west);
  \draw[-{Stealth}](z)-|(9.6,1.05)|-(b5.west);
  \node[itbaCyanDark] at (4.3,-1.6){$H^s_{k_2}(m)=z$};
  \node[itbaCyanDark] at (10.8,-1.6){$h^s_{k_1}(z)=t$};
\end{tikzpicture}
\caption{NMAC: hash con clave $k_2$ (y la longitud $L$ como bloque final) seguido de un paso
con $k_1$. \textit{(Recreación de la diapositiva práctica.)}}
\end{figure}

\subsection{HMAC (hash-based MAC)}
La construcción de HMAC está especificada en una RFC y es la más usada.

\begin{defbox}{HMAC}
\textbf{Gen:} emite $(s,k)$ con $|k|=n$.\quad
\[
  t := H^s_{\mathrm{IV}}\Bigl((k\xor\mathrm{opad})\concat
        H^s_{\mathrm{IV}}\bigl((k\xor\mathrm{ipad})\concat m\bigr)\Bigr),
\]
con $\mathrm{ipad}=\texttt{0x36}\cdots\texttt{36}$, $\mathrm{opad}=\texttt{0x5c}\cdots
\texttt{5c}$ (constantes de $n$ bytes) e $\mathrm{IV}$ una constante fija de longitud $n$.\\
\textbf{Es infalsificable (Mac-forge)} si $H$ es libre de colisiones.
\end{defbox}

\begin{slidebox}
\textbf{Eficiencia.} Aunque a primera vista parezca que se procesa el mensaje dos veces, en
realidad $m$ se procesa \textbf{una sola vez} en el hash interno; el paso externo (con
$\mathrm{opad}$) solo procesa la etiqueta intermedia $z$ (longitud fija). Por eso es muy
eficiente. El prefijo $k\xor\mathrm{ipad}$ \textbf{rompe el patrón} entre mensajes iguales y
los estados intermedios.
\end{slidebox}

% ---- Diagrama HMAC ----
\begin{figure}[h]
\centering
\begin{tikzpicture}[font=\scriptsize,>=Stealth,
   tr/.style={draw,trapezium,trapezium left angle=70,trapezium right angle=110,
              fill=itbaBlue!50,minimum width=1.5cm,minimum height=0.9cm,shape border rotate=180},
   blk/.style={draw,fill=slideGray,minimum width=1.7cm,minimum height=0.5cm},
   nout/.style={draw,fill=white,minimum width=0.8cm,minimum height=0.45cm}]
  % inner
  \node[blk](i0) at (0,1.4){$\mathrm{IV}\concat(k\xor\mathrm{ipad})$};
  \node[tr](hi0) at (0,0.4){$h^s$};\node[nout](z0) at (0,-0.65){$z_0$};
  \node[blk](i1) at (2.9,1.4){$z_0\concat m_1$};
  \node[tr](hi1) at (2.9,0.4){$h^s$};\node[nout](z1) at (2.9,-0.65){$z_1$};
  \node at (4.6,0.4){$\cdots$};
  \node[blk](i2) at (6.2,1.4){$z_{B}\concat L$};
  \node[tr](hi2) at (6.2,0.4){$h^s$};\node[nout,fill=itbaBlue!25](z) at (6.2,-0.65){$z$};
  \foreach \b/\h in {i0/hi0,i1/hi1,i2/hi2}{\draw[-{Stealth}](\b)--(\h);}
  \draw[-{Stealth}](hi0)--(z0);\draw[-{Stealth}](hi1)--(z1);\draw[-{Stealth}](hi2)--(z);
  \draw[-{Stealth}](z0)-|(1.45,1.05)|-(i1.west);
  \draw[-{Stealth}](z1)-|(4.6,1.05)--(4.6,1.05);
  \draw[-{Stealth}](4.85,1.05)-|(5.35,1.05)|-(i2.west);
  \node[itbaCyanDark,align=left] at (8.6,0.0){$H^s_{\mathrm{IV}}\bigl((k\xor\mathrm{ipad})\concat m\bigr)$\\$=z$};
  % outer
  \node[blk](o0) at (0,-2.4){$\mathrm{IV}\concat(k\xor\mathrm{opad})$};
  \node[tr](ho0) at (0,-3.4){$h^s$};\node[nout](w0) at (0,-4.45){$z_0$};
  \node[blk](o1) at (2.9,-2.4){$z_0\concat z$};
  \node[tr](ho1) at (2.9,-3.4){$h^s$};
  \node[draw,fill=practGreenL,minimum width=0.8cm](t) at (2.9,-4.45){$t$};
  \draw[-{Stealth}](o0)--(ho0);\draw[-{Stealth}](ho0)--(w0);
  \draw[-{Stealth}](o1)--(ho1);\draw[-{Stealth}](ho1)--(t);
  \draw[-{Stealth}](w0)-|(1.45,-2.75)|-(o1.west);
  \draw[-{Stealth},practGreen,thick](z.south)--(6.2,-1.65)-|(o1.north);
  \node[itbaCyanDark,align=left] at (6.6,-4.0){$H^s_{\mathrm{IV}}\bigl((k\xor\mathrm{opad})\concat z\bigr)=t$};
\end{tikzpicture}
\caption{HMAC: hash interno sobre $(k\xor\mathrm{ipad})\concat m$ produce $z$; hash externo
sobre $(k\xor\mathrm{opad})\concat z$ produce la etiqueta $t$.
\textit{(Recreación de la diapositiva práctica.)}}
\end{figure}

%=====================================================================
\section{Privacidad e integridad: combinarlas}
%=====================================================================

Para lograr \textbf{privacidad e integridad} a la vez hay que \emph{combinar} cifrado y
etiquetado. Hay \textbf{tres formas}, pero \textbf{solo dos son seguras}.

% ---- Diagrama semaforo de las 3 combinaciones ----
\begin{figure}[h]
\centering
\begin{tikzpicture}[font=\small,>=Stealth,node distance=4mm]
  \node[draw,rounded corners,fill=slideGray,minimum width=8.6cm,align=left] (c1)
    {\textbf{1. Cifrar y autenticar (por separado):}\quad
     $c\leftarrow\Enc_{k_1}(m),\ \ t\leftarrow\Mac_{k_2}(m)$\quad emite $\langle c,t\rangle$};
  \node[circle,fill=attackRed,minimum size=0.7cm,right=0.5cm of c1] (l1) {};
  \node[font=\scriptsize,right=1mm of l1,align=left] {\textbf{INSEGURO}\\ $t$ puede dar info de $m$};

  \node[draw,rounded corners,fill=slideGray,minimum width=8.6cm,align=left,below=4mm of c1] (c2)
    {\textbf{2. Autenticar, luego cifrar:}\quad
     $c\leftarrow\Enc_{k_1}\bigl(m\concat\Mac_{k_2}(m)\bigr)$};
  \node[circle,fill=yellow!80!orange,minimum size=0.7cm,right=0.5cm of c2] (l2) {};
  \node[font=\scriptsize,right=1mm of l2,align=left] {\textbf{Puede} ser seguro\\ requiere prueba (p.\,ej.\ \emph{padding})};

  \node[draw,rounded corners,fill=slideGray,minimum width=8.6cm,align=left,below=4mm of c2] (c3)
    {\textbf{3. Cifrar, luego autenticar:}\quad
     $c\leftarrow\Enc_{k_1}(m),\ \ t\leftarrow\Mac_{k_2}(c)$\quad emite $\langle c,t\rangle$};
  \node[circle,fill=practGreen,minimum size=0.7cm,right=0.5cm of c3] (l3) {};
  \node[font=\scriptsize,right=1mm of l3,align=left] {\textbf{SEGURO}\\ si ambos lo son y $k_1\neq k_2$};
\end{tikzpicture}
\caption{Las tres combinaciones de cifrado y autenticación, con su ``semáforo'' de
seguridad. \textit{(Recreación de las diapositivas teórica y práctica.)}}
\end{figure}

\begin{itemize}[leftmargin=1.4em,itemsep=2pt]
  \item \textbf{(1) Cifrar y autenticar} ($\Mac$ sobre el \emph{plano}): \textcolor{attackRed}{
        \textbf{inseguro}}. La etiqueta $t$ puede brindar información sobre $m$ (un MAC no
        garantiza confidencialidad de su entrada; dos mensajes iguales podrían dar la misma
        etiqueta).
  \item \textbf{(2) Autenticar luego cifrar} ($\Mac$ sobre el plano, todo cifrado):
        \textbf{puede ser seguro}, pero requiere prueba (problemas tipo \emph{padding}).
  \item \textbf{(3) Cifrar luego autenticar} ($\Mac$ sobre el \emph{cifrado}): \textcolor{
        practGreen}{\textbf{siempre seguro}}, siempre que los algoritmos de cifrado y MAC lo
        sean y las claves $k_1,k_2$ sean \textbf{independientes}. (Con claves iguales el
        esquema puede ser susceptible a ataques.)
\end{itemize}

%=====================================================================
\section{Cifrado Autenticado}
%=====================================================================

\begin{defbox}{Construcción de cifrado autenticado (encrypt-then-MAC)}
Sean $\Pi_e=(\Gen_e,\Enc,\Dec)$ un criptosistema \textbf{CPA-secure} y
$\Pi_m=(\Gen_m,\Mac,\Vrfy)$ un MAC \textbf{infalsificable} (con etiquetas únicas: para todo
$k,m$ hay un único $t$ con $\Vrfy_k(m,t)=1$). Se define:
\begin{align*}
  \textbf{Gen:}\ & k_1\leftarrow\Gen_e,\quad k_2\leftarrow\Gen_m \quad(\textbf{independientes})\\
  \textbf{Enc:}\ & c\leftarrow\Enc_{k_1}(m),\quad t\leftarrow\Mac_{k_2}(c) \quad\text{emitir }\langle c,t\rangle\\
  \textbf{Dec:}\ & \text{si } \Vrfy_{k_2}(c,t)=1 \to m=\Dec_{k_1}(c);\quad \text{si no } \to \perp\ (\text{fallo})
\end{align*}
\textbf{El criptosistema resultante es CCA-secure.} (Ver Teorema 4.20, Katz \& Lindell.) La
idea de la prueba es mostrar que el oráculo de descifrado \textbf{no le sirve de nada} al
adversario (cualquier $c$ adulterado falla la verificación del MAC).
\end{defbox}

% ---- Diagrama cifrado autenticado ----
\begin{figure}[h]
\centering
\begin{tikzpicture}[font=\small,>=Stealth,node distance=5mm,
   bx/.style={draw,rounded corners,minimum width=1.5cm,minimum height=0.8cm}]
  \node[bx,fill=slideGray] (m) at (0,0) {$m$};
  \node[bx,fill=itbaBlue] (enc) at (2.4,0) {$\Enc_{k_1}$};
  \node[bx,fill=slideGray] (c) at (4.8,0) {$c$};
  \node[bx,fill=practGreenL] (mac) at (7.2,0.0) {$\Mac_{k_2}$};
  \node[bx,fill=practGreenL] (t) at (7.2,-1.4) {$t$};
  \node[bx,fill=slideGray,minimum width=2.4cm] (pair) at (10.3,-0.7) {$\langle c,\,t\rangle$};
  \draw[-{Stealth}] (m)--(enc); \draw[-{Stealth}] (enc)--(c);
  \draw[-{Stealth}] (c)--(mac); \draw[-{Stealth}] (mac)--(t);
  \draw[-{Stealth}] (c.south) |- (pair.west);
  \draw[-{Stealth}] (t.east) -| (pair.south);
\end{tikzpicture}
\caption{Cifrado autenticado = \emph{encrypt-then-MAC}: privacidad ($\Enc$ CPA-secure) +
integridad ($\Mac$ infalsificable) $\Rightarrow$ \textbf{CCA-secure + MAC-forge secure}.}
\end{figure}

\begin{slidebox}
\textbf{Cifrado autenticado (Secure message transmission):}
\[
  \textbf{Privacidad} + \textbf{integridad} \;=\; \text{CCA-Secure} + \text{MAC-forge Secure}.
\]
\end{slidebox}

\subsection{Encadenamiento CCM}
\begin{defbox}{CCM — Counter with CBC-MAC}
Modo ``\textbf{Authenticate-then-encrypt}'':
\[
  e_k\bigl(\text{cbc-mac}_k(M)\concat m\bigr).
\]
Existe una \textbf{prueba de seguridad específica}: si el IV y el \emph{nonce} \textbf{no
coinciden ni se reutilizan}, la construcción es \textbf{CCA-secure usando LA MISMA CLAVE}.
\emph{(Ejercicio de la slide: esquematizar un cifrado usando AES-CCM.)}
\end{defbox}

\subsection{Encadenamiento GCM}
\begin{defbox}{GCM — Galois/Counter Mode}
Forma de encadenar un criptosistema de bloque que \textbf{provee cifrado autenticado}. El
cifrado es en \textbf{modo counter}; la etiqueta de autenticación (MAC) se calcula con
\textbf{GHASH} sobre $\mathrm{GF}(2^{128})$:
\[
  H=E_k(0^{128}),\qquad \mathrm{Mult}_h(x)=\mathrm{Mult}(x,h),\qquad
  \mathrm{Mult}(x,y)=x\cdot y \bmod \bigl(x^{128}+x^7+x^2+x+1\bigr).
\]
\end{defbox}

% ---- Diagrama GCM ----
\begin{figure}[h]
\centering
\begin{tikzpicture}[font=\scriptsize,>=Stealth,
   ctr/.style={draw,fill=white,minimum width=1.5cm,minimum height=0.5cm},
   ek/.style={draw,rounded corners,fill=itbaCyan,minimum width=1.1cm,minimum height=0.7cm},
   mh/.style={draw,rounded corners,fill=itbaCyan,minimum width=1.1cm,minimum height=0.6cm},
   xo/.style={circle,draw,inner sep=1.2pt},
   inc/.style={draw,rounded corners,fill=itbaCyan,minimum size=0.45cm}]
  % counters
  \node[ctr](c0) at (0,3.0){Counter 0};
  \node[inc](i0) at (1.6,3.0){incr};
  \node[ctr](c1) at (3.2,3.0){Counter 1};
  \node[inc](i1) at (4.8,3.0){incr};
  \node[ctr](c2) at (6.4,3.0){Counter 2};
  \draw[-{Stealth}](c0)--(i0);\draw[-{Stealth}](i0)--(c1);
  \draw[-{Stealth}](c1)--(i1);\draw[-{Stealth}](i1)--(c2);
  % E_k
  \node[ek](e0) at (0,2.0){$E_k$};\node[ek](e1) at (3.2,2.0){$E_k$};\node[ek](e2) at (6.4,2.0){$E_k$};
  \draw[-{Stealth}](c0)--(e0);\draw[-{Stealth}](c1)--(e1);\draw[-{Stealth}](c2)--(e2);
  % plaintext xor
  \node[ctr](p1) at (2.0,1.1){Plaintext 1};
  \node[xo](xp1) at (3.2,1.1){$\xor$};
  \node[ctr](p2) at (5.2,1.1){Plaintext 2};
  \node[xo](xp2) at (6.4,1.1){$\xor$};
  \draw[-{Stealth}](e1)--(xp1);\draw[-{Stealth}](p1)--(xp1);
  \draw[-{Stealth}](e2)--(xp2);\draw[-{Stealth}](p2)--(xp2);
  % ciphertext
  \node[ctr](ct1) at (3.2,0.2){Ciphertext 1};
  \node[ctr](ct2) at (6.4,0.2){Ciphertext 2};
  \draw[-{Stealth}](xp1)--(ct1);\draw[-{Stealth}](xp2)--(ct2);
  % GHASH chain
  \node[mh](m0) at (1.2,-0.9){$\mathrm{mult}_H$};
  \node[xo](g0) at (1.2,0.2){$\xor$};
  \node[mh](m1) at (3.2,-1.6){$\mathrm{mult}_H$};
  \node[xo](g1) at (3.2,-0.9){$\xor$};
  \node[mh](m2) at (6.4,-0.9){$\mathrm{mult}_H$};
  \node[xo](g2) at (6.4,-0.2){$\xor$};
  \node[ctr](ad) at (-0.4,-0.9){Auth Data 1};
  \node[ctr](len) at (5.0,-1.6){len(A)$\|$len(C)};
  \draw[-{Stealth}](ad)--(m0);
  \draw[-{Stealth}](g0)--(m0.north);
  \draw[-{Stealth}](m0)|-(2.4,-1.6)--(m1.west);
  \draw[-{Stealth}](ct1)--(g1);\draw[-{Stealth}](g1)--(m1);
  \draw[-{Stealth}](ct2)--(g2);\draw[-{Stealth}](g2)--(m2);
  \draw[-{Stealth}](m1)|-(4.4,-2.1)-|(g2.south);
  \draw[-{Stealth}](len)--(m2.west);
  \node[ctr,fill=practGreenL](at) at (6.4,-2.4){Auth Tag};
  \draw[-{Stealth}](m2)--(at);
\end{tikzpicture}
\caption{GCM: cifrado en modo counter más cadena GHASH (multiplicaciones en
$\mathrm{GF}(2^{128})$) que produce el \emph{Auth Tag}. \textit{(Recreación de la
diapositiva.)}}
\end{figure}

\begin{practbox}{Ejercicio (slide práctica): claves iguales rompen encrypt-then-MAC}
\textbf{Enunciado.} Considerar $F$ una función pseudoaleatoria biyectiva fuerte,
$\Enc_k(m)=F_k(m\concat r)$ y $\Mac_k(c)=F^{-1}_k(c)$ (pasa Mac-forge). Probar que el esquema
\emph{cifrar y luego autenticar} \textbf{no es resistente a CCA por tener claves iguales}.

\medskip
\textbf{Idea de la resolución.} Como $\Mac_k=F^{-1}_k$ y $\Enc_k=F_k(\cdot)$ usan
\emph{la misma} clave $k$, el tag $t=F^{-1}_k(c)=F^{-1}_k(F_k(m\concat r))=m\concat r$
\textbf{revela el plano} (más el random). El adversario, al recibir el challenge
$\langle c,t\rangle$, lee $t$ y obtiene directamente $m\concat r$, distinguiendo $m_0$ de
$m_1$ \emph{sin} siquiera usar el oráculo de descifrado. La moraleja confirma la condición de
la slide: \textbf{las claves $k_1$ y $k_2$ deben ser independientes}.
\end{practbox}

%=====================================================================
\section{Aplicaciones prácticas}
%=====================================================================

\begin{slidebox}
\begin{itemize}[leftmargin=1.3em,itemsep=2pt]
  \item En términos generales, \textbf{usar solo cifrado autenticado} en sistemas reales.
  \item El cifrado ``normal'' puede ser \textbf{manipulado} (maleabilidad).
  \item En aplicaciones reales, el requerimiento de \textbf{privacidad lleva implícito el de
        integridad}.
  \item Muy pocas veces se tiene control sobre el material que va a ser cifrado y descifrado:
        la solución debe funcionar \textbf{independientemente} de eso.
\end{itemize}
\end{slidebox}

%=====================================================================
\section*{Cross-check teoría vs.\ práctica}
\addcontentsline{toc}{section}{Cross-check teoría vs.\ práctica}
%=====================================================================
\begin{practbox}{Qué agrega la clase práctica}
\begin{itemize}[leftmargin=1.4em,itemsep=3pt]
  \item \textbf{Mac-forge explícito sobre los 3 MACs inseguros:} la práctica (y la teórica)
        resuelven \emph{paso a paso} cómo se rompe $G(k)\xor m$, $k\xor\mathrm{first\_k\_bits}
        (m)$ y $\Enc_k(|m|)$, destilando las propiedades de un MAC seguro (procesar todo el
        mensaje, sin determinismos extraíbles).
  \item \textbf{Análisis de las 3 sugerencias para longitud variable} (XOR de bloques, tag
        por bloque, tag con número de secuencia) con el ataque concreto de cada una
        (anular el XOR, permutar bloques, mezclar bloques de mensajes distintos en la misma
        posición). Esto es desarrollo \emph{propio de la práctica}, no de las slides
        teóricas.
  \item \textbf{Construcción segura $r\concat L\concat i\concat m_i$} y la justificación
        intuitiva de por qué el random $r$ impide mezclar mensajes (no hay un único $r$ que
        sirva para un tag híbrido), más la observación de que es \textbf{ineficiente}.
  \item \textbf{Ataque a CBC-MAC de longitud variable} explicado vía estados intermedios
        ($t_1,t_2$ que se descartan pero se recalculan).
  \item \textbf{Aclaraciones del profesor} ausentes en las slides: el MAC \textbf{no provee
        no-repudio} (clave simétrica) y \textbf{no protege contra replay} (usar número de
        secuencia o timestamp).
  \item \textbf{Ejercicio CCA con claves iguales} ($\Enc_k(m)=F_k(m\concat r)$,
        $\Mac_k(c)=F^{-1}_k(c)$) que motiva la independencia de claves.
\end{itemize}
\end{practbox}

%=====================================================================
\section*{Lectura recomendada}
\addcontentsline{toc}{section}{Lectura recomendada}
%=====================================================================
\begin{center}
\begin{tcolorbox}[enhanced,colback=itbaBlue!20,colframe=itbaCyanDark,boxrule=1pt,
  arc=3pt,width=0.85\textwidth,halign=center]
\large\textbf{Capítulo 4}\\[4pt]
\emph{Introduction to Modern Cryptography}\\[2pt]
Katz \& Lindell
\end{tcolorbox}
\end{center}

\noindent\textbf{Guías de práctica:} se sugiere avanzar hasta la \textbf{guía 3}.

%---------------------------------------------------------------------
\vfill
\begin{center}\small\itshape\color{black!60}
Apunte integral de la Clase 3 — MACs y Cifrado Autenticado · Criptografía y Seguridad, ITBA.\\
Síntesis de teoría (transcripciones + diapositivas) y práctica.
\end{center}

\end{document}
